%% General-TBox ALC: type elimination and combined complexity.

\section{General-TBox \(\mathcal{ALC}\)}

\begin{annexitem}{annex:p9.alc.2}{thm:p09-alc-type-elimination}{Correctness and bounds of \(\mathcal{ALC}\) type elimination}
\begin{proof}
Write \(\mathcal C\) for the finite closure and \(S_i\) for the
successive type sets in Method~\ref{def:p09-alc-type-elimination}.
The NNF complement is an involution and represents semantic
complementation, by the recursive Boolean and modal truth clauses.

\paragraph{Types realized in a model survive.}
Suppose \(I\) satisfies the input, including \(C_0^I\ne\varnothing\)
when a concept is tested.
For \(x\in\Delta^I\), put
\[
 \operatorname{tp}_I(x)=\{E\in\mathcal C:x\in E^I\}.
\]
It belongs to \(S_0\): the constant, complement and Boolean conditions
are truth clauses, and each internalized GCI is true at every
domain element.
We prove simultaneously for all \(x\) that
\(\operatorname{tp}_I(x)\in S_i\) at every stage \(i\).
If this holds at \(i\) and \(\exists r.E\in\operatorname{tp}_I(x)\),
choose \(y\) with \((x,y)\in r^I\) and \(y\in E^I\).
Its type is in \(S_i\) by the induction hypothesis.
Every \(\forall r.F\) true at \(x\) has \(F\) true at \(y\),
so
\[
 \operatorname{tp}_I(x)\leadsto_r\operatorname{tp}_I(y),
 \qquad E\in\operatorname{tp}_I(y).
\]
The existential requirement therefore has a witness, and the type
of \(x\) is retained at stage \(i+1\).

At stabilization, all these realized types belong to \(S\).
The domain is nonempty, so \(S\ne\varnothing\).
Assign \(t_a=\operatorname{tp}_I(a^I)\) to each ABox name.
Concept assertions and role assertions satisfy the two acceptance
conditions by the same argument.
If \(C_0\) is tested, its chosen semantic witness supplies a type
in \(S\) containing \(\operatorname{nnf}(C_0)\).
Thus every satisfiable input is accepted.

\paragraph{The canonical interpretation.}
Conversely, suppose the method accepts, with a witnessing assignment
\(a\mapsto t_a\).
Let \(I_*\) have domain \(S\), and interpret the input names by
\[
\begin{aligned}
 A^{I_*}&=\{t\in S:A\in t\}
              &&\text{for concept names occurring in the closure},\\
 r^{I_*}&=\{(t,u)\in S^2:t\leadsto_r u\}
              &&\text{for input role names},\\
 a^{I_*}&=t_a
              &&\text{for ABox individual names}.
\end{aligned}
\]
Unused concept names and roles can be empty, and unused individuals
can denote any fixed member of the nonempty set \(S\).
No injectivity requirement is imposed on the individual assignment.

For every \(E\in\mathcal C\) and \(t\in S\), we prove
\[
 t\in E^{I_*}\quad\Longleftrightarrow\quad E\in t.
\]
Induct on NNF rank, counting concept-name literals, \(\top\), and
\(\bot\) as rank zero.
The definition of the interpretation proves the positive-literal
case.  Exactly one of \(A,\neg A\) belongs to a type, proving the
negative-literal case.  The constant and Boolean cases follow from
the corresponding type conditions.

If \(\exists r.F\in t\), stability of \(S\) supplies a compatible
\(u\in S\) containing \(F\); induction gives the required successor.
If \(\exists r.F\notin t\), completeness of the type puts
\(\forall r.\bar F\) in \(t\).
Every compatible successor then contains \(\bar F\), so by induction
and semantic complementation none satisfies \(F\).
Thus the existential concept is false.

If \(\forall r.F\in t\), every compatible successor contains \(F\)
and satisfies it by induction, so the universal is true.
If \(\forall r.F\notin t\), its complement
\(\exists r.\bar F\) belongs to \(t\).
Stability supplies a compatible successor containing \(\bar F\);
induction makes this a counterexample to the universal.
Both \(F\) and \(\bar F\) have lower rank than the modal concept,
so none of these cases invokes the desired equivalence circularly.
This completes the truth lemma.

Every member of \(\mathsf{Glob}(\mathcal T)\) belongs to every type,
so the truth lemma proves all TBox inclusions globally.
The selected \(t_a\) satisfy the ABox concepts, and each asserted
role edge is one of the compatible pairs used in the interpretation.
The tested concept, when present, belongs to some surviving type
and is therefore nonempty in this same model.
This proves soundness of acceptance.
Its domain has size \(|S|\leq2^m\).

\paragraph{Termination and deterministic time.}
NNF conversion visits each constructor of the input concepts once.
Taking sub-concepts and their NNF complements yields \(m=O(N)\)
distinct closure expressions, each with a representation of
polynomial length.  Equivalently, use shared pointers for the
subexpressions and one complemented copy of each.
The number \(M\) of initial types is at most \(2^m\).
Enumerating all subsets and checking the local conditions costs
\(2^m\operatorname{poly}(N)\).

An elementary implementation scans every remaining type, every one
of its at most \(m\) existential requirements, and all candidate
witness types.  Compatibility is a polynomial-time test on their
bit-vector descriptions.
At most \(M\) scans make a strict deletion before stabilization.
Thus even this implementation takes at most
\(M^3\operatorname{poly}(N)\) time.
If \(k\) distinct individual names occur, enumerate at most \(M^k\)
assignments and check each finite list of assertions.
The optional-concept and nonempty-domain tests take polynomial
time in \(M+N\).
Overall the bound is
\[
 \operatorname{poly}(N)\bigl(2^{3m}+2^{mk}\bigr)
       \leq 2^{O(N^2)},
\]
since \(m,k=O(N)\).
For \(k=0\), the bound is \(2^{O(N)}\).
In particular, the wholly empty input leaves the type containing
\(\top\), and the nonempty-domain test accepts it.
A TBox forcing \(\top\sqsubseteq\bot\) leaves no type and is rejected.
\end{proof}
\end{annexitem}

\begin{annexitem}{annex:p9.alc.3}{thm:p09-alc-general-exptime}{Combined complexity of general-TBox \(\mathcal{ALC}\)}
\begin{proof}
The type-elimination bound in Theorem~\ref{thm:p09-alc-type-elimination}
places knowledge-base satisfiability, with or without a tested concept,
in \(\mathsf{EXPTIME}\).
We prove hardness first for a tested concept with a general TBox and
an empty ABox.

\paragraph{The alternating-machine source problem.}
Use the characterization
\[
 \mathsf{APSPACE}=\mathsf{EXPTIME}
\]
from \cite[Corollary~3.6, p.~122]{ChandraKozenStockmeyer1981}.
Here alternating polynomial space uses existential and universal
computation states.
Fix an arbitrary language in \(\mathsf{EXPTIME}\), and choose one
fixed alternating machine \(M\) for it, with a polynomial space bound.
The source's space definition initially bounds an accepting computation
tree.  Apply its Theorem~2.5 and Theorem~2.6(b), using a
space-constructible polynomial bound, to choose an equivalent machine
without negating states whose paths all terminate and whose reachable
configurations all obey that bound
\cite[pp.~120--122]{ChandraKozenStockmeyer1981}.
On input \(w\), use an equivalent one-tape form whose tape occupies
cells \(1,\ldots,s\), where \(s=s(|w|)\geq\max(1,|w|)\) is polynomial.
A fixed number of tapes can be encoded by tracks and marked head
positions on one tape; finite scans collect the symbols for a step
and perform the selected update.  This uses polynomial space.

Let the finite control set be the disjoint union
\[
 Q=Q_\exists\mathbin{\dot\cup}Q_\forall
                   \mathbin{\dot\cup}\{q_{\rm acc},q_{\rm rej}\},
\]
and let \(\Gamma\) be the finite tape alphabet.
Normalize every nonhalting situation to two instructions
\[
 \delta_j(q,\gamma)=(q',\gamma',d),
 \qquad j\in\{0,1\},\quad d\in\{-1,0,1\}.
\]
Single or deterministic choices are duplicated.  Larger finite
choices are replaced by binary trees of the same existential or
universal kind.
If an initial machine convention admits a nonhalting situation with
no moves, replace an empty existential choice by rejection and an
empty universal choice by acceptance, and then use two identical
moves to that halting answer.
The two designated halting controls have no instructions.
These normalizations, together with removal of negating states,
are compatible with the cited alternation characterization
\cite[\S2]{ChandraKozenStockmeyer1981}.

A bounded configuration is \(c=(q,i,u)\), with \(q\in Q\),
\(1\leq i\leq s\), and \(u\in\Gamma^s\).
If an instruction would leave the tape interval, its successor
instead has rejecting control, the same valid head position, and
the specified tape write.
This gives two well-defined successors \(\tau_0(c),\tau_1(c)\)
of every nonhalting bounded configuration.
The chosen space bound makes this convention agree with the machine
on its actual computations from the input.

An accepting computation is a finite tree of configurations.
Its leaves have control \(q_{\rm acc}\).
At an existential node it has one child for one of the two
instructions; at a universal node it has both indexed children.
No rejecting configuration occurs.
Two equal successor configurations may appear as distinct tree nodes.

\paragraph{A finite bound on an accepting tree.}
There are at most
\[
 K=|Q|\,s\,|\Gamma|^s
\]
bounded configurations.
Starting with \(W_0=\varnothing\), define \(W_{i+1}\) to contain
all accepting configurations, the existential configurations with
at least one successor in \(W_i\), and the universal configurations
with both successors in \(W_i\).
Induction gives that \(W_{i+1}\) is exactly the set of configurations
having an accepting tree of height at most \(i\), with leaves of
height zero.
The sequence is increasing, and an equality at one stage persists
thereafter.  It stabilizes by \(W_K\), because its underlying
configuration set has at most \(K\) elements.
Thus, whenever the machine accepts, it has an accepting tree of
height at most \(K-1\).

\paragraph{Concept names for configurations.}
Use role names \(r_0,r_1\) and concept names
\[
 \mathsf Q_q\ (q\in Q),\qquad
 \mathsf H_i\ (1\leq i\leq s),\qquad
 \mathsf A_{i,\gamma}\ (1\leq i\leq s,\ \gamma\in\Gamma).
\]
They record the control, head position, and tape symbols.
Introduce \(b\) rank-bit names \(\mathsf B_0,\ldots,\mathsf B_{b-1}\),
with bit zero least significant, where
\[
 b=1+\lceil\log_2|Q|\rceil+\lceil\log_2s\rceil
                                  +s\lceil\log_2|\Gamma|\rceil.
\]
Then \(2^{b-1}\geq K\), so \(R=2^b-1\geq K\).
Only the \(b\) bit names are written in the output, not a unary
representation of \(R\).

The TBox first requires exactly one control and one head position:
\[
\begin{gathered}
 \top\sqsubseteq\bigsqcup_{q\in Q}\mathsf Q_q,\qquad
 \mathsf Q_q\sqcap\mathsf Q_{q'}\sqsubseteq\bot\quad(q\ne q'),\\
 \top\sqsubseteq\bigsqcup_{i=1}^{s}\mathsf H_i,\qquad
 \mathsf H_i\sqcap\mathsf H_{i'}\sqsubseteq\bot\quad(i\ne i').
\end{gathered}
\]
For each cell \(i\), require exactly one tape symbol:
\[
 \top\sqsubseteq\bigsqcup_{\gamma\in\Gamma}\mathsf A_{i,\gamma},
 \qquad
 \mathsf A_{i,\gamma}\sqcap\mathsf A_{i,\eta}\sqsubseteq\bot
                                      \quad(\gamma\ne\eta).
\]
Each model element therefore describes one bounded configuration.
Different elements may describe the same one.
For halting and branching, include
\[
\begin{aligned}
 \mathsf Q_{q_{\rm rej}}&\sqsubseteq\bot,\\
 \mathsf Q_{q_{\rm acc}}&\sqsubseteq
                   \forall r_0.\bot\sqcap\forall r_1.\bot,\\
 \mathsf Q_q&\sqsubseteq
                   \exists r_0.\top\sqcup\exists r_1.\top
                         &&(q\in Q_\exists),\\
 \mathsf Q_q&\sqsubseteq
                   \exists r_0.\top\sqcap\exists r_1.\top
                         &&(q\in Q_\forall).
\end{aligned}
\]

\paragraph{Updating one cell and preserving all others.}
For every nonhalting \(q\), head position \(i\), symbol \(\gamma\),
and instruction index \(j\), let
\[
 G_{q,i,\gamma}
       =\mathsf Q_q\sqcap\mathsf H_i\sqcap\mathsf A_{i,\gamma}.
\]
Compute the instruction's bounded target data \((q',i',\gamma')\),
including the rejecting boundary convention above, and add
\[
 G_{q,i,\gamma}\sqsubseteq
       \forall r_j.(\mathsf Q_{q'}\sqcap\mathsf H_{i'}
                                      \sqcap\mathsf A_{i,\gamma'}).
\]
For every \(i,\gamma,j\), also add the unchanged-cell condition
\[
 \neg\mathsf H_i\sqcap\mathsf A_{i,\gamma}
                     \sqsubseteq\forall r_j.\mathsf A_{i,\gamma}.
\]
Exactly one guard \(G_{q,i,\gamma}\) is true at a nonhalting element.
Its update determines the new control, head and scanned cell.
The unchanged-cell conditions determine every other cell.
Consequently every \(r_j\)-successor describes exactly
\(\tau_j(c)\), where \(c\) is the source configuration.
A rejecting target is impossible because rejecting control was
forbidden.  None of these conditions requires a role to be functional.

\paragraph{Forcing a strict binary countdown.}
Use the concept abbreviations
\[
 Z=\bigsqcap_{\ell<b}\neg\mathsf B_\ell,\qquad
 L_\ell=\bigsqcap_{h<\ell}\neg\mathsf B_h,
 \qquad L_0=\top.
\]
For each \(\ell<b\) and each \(r\in\{r_0,r_1\}\), include all four GCIs
\[
\begin{aligned}
 L_\ell\sqcap\mathsf B_\ell
    &\sqsubseteq\forall r.\neg\mathsf B_\ell,\\
 L_\ell\sqcap\neg\mathsf B_\ell
    &\sqsubseteq\forall r.\mathsf B_\ell,\\
 \neg L_\ell\sqcap\mathsf B_\ell
    &\sqsubseteq\forall r.\mathsf B_\ell,\\
 \neg L_\ell\sqcap\neg\mathsf B_\ell
    &\sqsubseteq\forall r.\neg\mathsf B_\ell.
\end{aligned}
\]
Also include
\[
 Z\sqsubseteq\forall r_0.\bot\sqcap\forall r_1.\bot.
\]
At an interpretation element \(x\), define the integer
\[
 \rho(x)=\sum_{\ell<b}2^\ell\,
                    \mathbf 1_{x\in\mathsf B_\ell^I}.
\]
If \(\rho(x)>0\), let \(\ell_0\) be its least set bit.
For \(\ell<\ell_0\), all lower bits and bit \(\ell\) are zero,
so the second GCI makes that successor bit one.
At \(\ell_0\), the first GCI changes the bit from one to zero.
For \(\ell>\ell_0\), a lower bit is already one at the source;
the third and fourth GCIs preserve the source bit.
These are exactly the bits of \(\rho(x)-1\).
At rank zero, the final GCI forbids both roles.
Thus every role edge in every model strictly decreases this
nonnegative integer, with no wraparound.

\paragraph{The initial concept and polynomial output size.}
Let \((q_{\rm init},1,u_1\cdots u_s)\) be the initial configuration
for \(w\), with its input followed by blanks in the chosen convention.
Use the tested concept
\[
 C_{M,w}=\mathsf Q_{q_{\rm init}}\sqcap\mathsf H_1
            \sqcap\bigsqcap_{i=1}^{s}\mathsf A_{i,u_i}
            \sqcap\bigsqcap_{\ell<b}\mathsf B_\ell.
\]
Let \(\mathcal T_{M,s}\) be exactly the finite GCI families displayed
above.  The symbols \(G_{q,i,\gamma},Z,L_\ell\) are abbreviations
for concepts and are expanded when measuring size.
The size is bounded by
\[
 O\!\left(|Q|^2+s^2+s|\Gamma|^2
                   +|Q|s|\Gamma|+b^2\right)
\]
before the logarithmic encoding of indices.
The \(b^2\) term includes all expanded lower-bit conjunctions.
For the fixed source machine, \(s\) and \(b\) are polynomial in
\(|w|\), so the entire TBox and initial concept are constructed
in polynomial time.

\paragraph{An accepting computation gives a model.}
If \(M\) accepts \(w\), choose an accepting tree of height at most
\(K-1\).  Take its nodes as the interpretation domain.
Interpret the configuration names by their node labels.
At an existential node include only its selected \(r_j\)-edge;
at a universal node include the two indexed edges to its two children.
An accepting leaf has neither edge.
Give a node at depth \(d\) the binary value \(R-d\).
This is nonnegative because the height is below \(K\leq R\).

The configuration partitions and all branching/halting GCIs hold
by the definition of the tree.
The instruction and unchanged-cell GCIs hold because its children
are the actual successors.
The rank decreases by exactly one on every edge, so the decrement
GCIs hold, and a rank-zero node, if present, is a leaf.
The root satisfies \(C_{M,w}\).
Thus \(C_{M,w}\) is satisfiable relative to \(\mathcal T_{M,s}\).

\paragraph{A model gives an accepting computation.}
Conversely, let \(I\models\mathcal T_{M,s}\) and let
\(x_0\in C_{M,w}^I\).
Every element describes exactly one bounded configuration, and no
element has rejecting control.
We prove by induction on \(\rho(x)\) that the configuration described
by each \(x\) has a finite accepting tree.
At rank zero there are no outgoing role edges, so the branching
GCIs exclude both nonhalting control classes.  The control must
be \(q_{\rm acc}\), giving an accepting leaf.
At positive rank, an accepting control again gives a leaf.
For an existential control choose an existing successor under one
of the two roles.  Its configuration is the corresponding machine
successor, and its rank is one less, so the induction hypothesis
supplies its accepting subtree.
For a universal control choose one successor for each role and
attach the two accepting subtrees supplied by induction.
These choices are finite at each stage, and the decreasing integer
rank bounds the height of the resulting finite tree.
The tree at \(x_0\) starts with the actual input configuration,
so \(M\) accepts \(w\).

We have proved
\[
 w\in L(M)\quad\Longleftrightarrow\quad
 C_{M,w}\text{ is satisfiable relative to }\mathcal T_{M,s}.
\]
Since the source language was arbitrary in \(\mathsf{EXPTIME}\),
this proves the required hardness using plain \(\mathcal{ALC}\).

\paragraph{The other reasoning tasks.}
For a fresh individual \(a_*\),
\[
 C\text{ is satisfiable relative to }\mathcal T
 \quad\Longleftrightarrow\quad
 (\mathcal T,\{C(a_*)\})\text{ is satisfiable}.
\]
This polynomial reduction proves knowledge-base hardness, and
the optional-tested-concept task already contains concept
satisfiability as its empty-ABox case.

Subsumption satisfies
\[
 C\sqsubseteq_{\mathcal T}D
 \quad\Longleftrightarrow\quad
 C\sqcap\neg D\text{ is unsatisfiable relative to }\mathcal T.
\]
Deterministic exponential time is closed under complement.
The complement of an \(\mathsf{EXPTIME}\)-complete language is
also \(\mathsf{EXPTIME}\)-hard: reduce the complement of an arbitrary
\(\mathsf{EXPTIME}\) language to the original complete language,
then complement the equivalence.
Thus concept unsatisfiability is \(\mathsf{EXPTIME}\)-complete,
and its reduction to \(C\sqsubseteq_{\mathcal T}\bot\) proves
subsumption hardness.  The displayed satisfiability reduction
also proves subsumption membership.

Finally,
\[
 (\mathcal T,\mathcal A)\models D(a)
 \quad\Longleftrightarrow\quad
 (\mathcal T,\mathcal A\cup\{\neg D(a)\})
                         \text{ is unsatisfiable},
\]
where the assertion is added to the ABox.
This gives the exponential upper bound for instance checking.
Subsumption reduces to instance checking by taking a fresh \(a_*\),
the knowledge base \((\mathcal T,\{C(a_*)\})\), and the query
\(D(a_*)\).
Every counterexample element to subsumption can interpret the fresh
name, and every countermodel to that instance query is such a
counterexample.  This proves hardness, including the vacuous cases
when the concept or knowledge base is unsatisfiable.
All the reductions are polynomial and stay within the stated input syntax.
\end{proof}
\end{annexitem}
